\documentclass[10pt]{article}
\usepackage[margin=1in]{geometry}
\usepackage[T1]{fontenc}
\usepackage{lmodern,microtype,amsmath,amssymb,amsthm}
\theoremstyle{definition}
\newtheorem{assumption}{Assumption}
\begin{document}
\section*{Scalar H\"older Stochastic Stationarity}
We seek a stationary point of a differentiable, possibly nonconvex function $f:\mathbb{R}\to\mathbb{R}$ whose derivative is H\"older continuous with constant $H>0$ and exponent $\nu\in(0,1]$. From $x^0\in\mathbb{R}$, the objective gap to some lower bound is at most $\Delta>0$. Each scalar oracle reply $g(x,\omega)$ at a query $x\in\mathbb{R}$ uses an independent sample $\omega$ from a distribution $\mu$ and is unbiased for $f'(x)$, with centered $p$th moment at most $\sigma^p$, where $p\in(1,2]$ and $\sigma\geq0$. The required output $\bar{x}\in\mathbb{R}$ satisfies $\mathbb{P}(|f'(\bar{x})|>\varepsilon)\leq 1/3$ for $\varepsilon>0$, where $\mathbb{P}$ is the joint law of algorithm and oracle randomness.

The task can be written as
\[
\inf_{x\in\mathbb{R}} f(x),
\]
with success measured by the derivative at the returned point. The measurable sample space $\Omega$ and the tuple $(H,\nu,\Delta,\sigma,p,\varepsilon,x^0)$ are public inputs. The objective $f$, sample distribution $\mu$, and oracle map $g$ specify an admissible instance.

\begin{assumption}[H\"older derivative]
The function $f$ is differentiable at every $x\in\mathbb{R}$, and
\[
|f'(x)-f'(y)|\leq H|x-y|^{\nu}
\qquad\text{for all }x,y\in\mathbb{R}.
\]
\end{assumption}

\begin{assumption}[Initial gap to a lower bound]
There is a number $\ell\in\mathbb{R}$ such that
\[
\ell\leq f(x)\quad\text{for all }x\in\mathbb{R},
\qquad f(x^0)-\ell\leq\Delta.
\]
\end{assumption}

\begin{assumption}[Scalar stochastic-gradient oracle]
There is a probability measure $\mu$ on $\Omega$ and a jointly measurable map $g:\mathbb{R}\times\Omega\to\mathbb{R}$. For every fixed $x\in\mathbb{R}$, both $g(x,\cdot)$ and $|g(x,\cdot)-f'(x)|^p$ are integrable under $\mu$, and
\begin{align*}
\int_{\Omega}g(x,\omega)\,d\mu(\omega)&=f'(x),\\
\int_{\Omega}|g(x,\omega)-f'(x)|^p\,d\mu(\omega)&\leq\sigma^p.
\end{align*}
The oracle samples $\omega_0,\omega_1,\ldots$ have a joint law under which the coordinate maps are measurable and independent, each with marginal $\mu$. Its reply at call $k$ and query $x$ is $g(x,\omega_k)\in\mathbb{R}$.
\end{assumption}

Queries may be chosen adaptively anywhere in $\mathbb{R}$. Each reply costs one oracle call, and any internal algorithm randomness is independent of the oracle samples. The algorithm is selected from the public inputs; an externally specified budget $K\in\mathbb{N}$ may depend on the full public tuple.

After $K$ calls, the algorithm returns a measurable scalar output $\bar{x}$. Under the joint law $\mathbb{P}$ of its randomness and the oracle samples, the accuracy requirement is
\[
\mathbb{P}\bigl(|f'(\bar{x})|>\varepsilon\bigr)\leq\frac{1}{3}.
\]
\end{document}
